% =====================================================================
%  So sánh: CORING gốc vs Rebuild YOLO vs Rebuild ResNet-56
%  Dự án: Component-aware Bi-level Pruning for One-stage Detectors
%
%  *** BIÊN DỊCH BẰNG XeLaTeX HOẶC LuaLaTeX *** (KHÔNG dùng pdfLaTeX).
%    Overleaf: Menu -> Compiler -> XeLaTeX.
% =====================================================================
\documentclass[11pt,a4paper]{article}

\usepackage{fontspec}
\setmainfont{Latin Modern Roman}     % font Unicode có đủ dấu tiếng Việt

\usepackage[margin=2cm]{geometry}
\usepackage{array}
\usepackage{tabularx}
\usepackage{enumitem}
\usepackage{hyperref}
\hypersetup{colorlinks=true, linkcolor=blue!50!black, urlcolor=blue!50!black}

\setlist{leftmargin=1.4em,topsep=2pt,itemsep=1pt}
\newcommand{\code}[1]{\texttt{#1}}
\newcolumntype{Y}{>{\raggedright\arraybackslash}X}

\title{\textbf{So sánh CORING: bản gốc vs 2 bản rebuild}\\
\large Bản gốc \textperiodcentered{} Rebuild YOLO (so bi-level) \textperiodcentered{} Rebuild ResNet-56 (validate)}
\author{}
\date{}

\begin{document}
\maketitle
\vspace{-2em}

% ---------------------------------------------------------------------
\section{Ba bản, ba mục đích}
\begin{itemize}
  \item \textbf{CORING gốc} (\href{https://github.com/pvti/CORING}{pvti/CORING}): phương pháp gốc
        --- \code{main.py} (one-shot) và \code{main\_kshot.py} (k-shot).
  \item \textbf{Rebuild YOLO} (\code{pruning/coring.py} + \code{prune.py} \code{-{}-method coring}):
        cắm CORING vào harness YOLO để \emph{so công bằng với bi-level} trên DeepFish.
  \item \textbf{Rebuild ResNet-56} (\code{prune\_resnet56\_cifar10.py}): chạy lại CORING trên chính
        benchmark gốc (ResNet-56/CIFAR-10) để \emph{kiểm chứng} thuật toán đúng (paper Table 3:
        94.76\% @ 22.4\% params).
\end{itemize}

% ---------------------------------------------------------------------
\section{Bảng so sánh}
\footnotesize
\renewcommand{\arraystretch}{1.3}
\noindent\makebox[\textwidth][c]{%   % bảng tràn ra 2 lề, không cần package
\begin{tabularx}{1.2\textwidth}{|p{2cm}|Y|Y|Y|p{2.5cm}|}
\hline
\textbf{Chiều} & \textbf{CORING gốc} & \textbf{Rebuild YOLO} & \textbf{Rebuild ResNet-56} & \textbf{Giống gốc?}\\
\hline
Mục đích & phương pháp gốc & so với bi-level (YOLO/DeepFish) & validate reproduce số paper (CIFAR-10) & Khác (chủ đích)\\
\hline
Thuật toán chọn & \code{get\_rank} (HOSVD $+$ VBD $+$ min\_sum, loop) & \code{get\_rank} (HOSVD $+$ VBD $+$ min\_sum, loop) & \code{get\_rank} (HOSVD $+$ VBD $+$ min\_sum, loop) & \textbf{Giống} (cả 2)\\
\hline
Get-prunable / surgery & \code{load\_*\_model} per-arch (copy top-K conv$+$fc) & \textbf{PrunableConv$+$mask$+$\code{convert\_to\_lean}} (harness mình) & \code{load\_resnet56\_pruned} (port từ gốc) & R56 giống \textperiodcentered{} YOLO khác\\
\hline
Copy BN sliced & \textbf{KHÔNG} (BN được init weight) & CÓ (copy như Conv $+$ Act) & CÓ (copy như Conv $+$ Act) & \textbf{Khác} (cả 2 thêm BN)\\
\hline
Schedule & \textbf{back-loaded} $r/(K{-}\text{shot})$ (\code{main\_kshot.py}) & \textbf{back-loaded} $r/(K{-}it)$ (\code{prune.py:155}) & \textbf{back-loaded} $r/(K{-}\text{shot})$ (\code{:278}) & \textbf{Giống} (cả 2)\\
\hline
Phân bổ per-layer & \textbf{LIST per-layer} (vd \code{[0.]+[0.18]*29}) & \textbf{UNIFORM} (1 scalar mọi layer) & \textbf{LIST per-layer} (như gốc) & R56 giống \textperiodcentered{} YOLO khác\\
\hline
Re-rank mỗi shot & có & có & có & \textbf{Giống} (cả 2)\\
\hline
Finetune & MultiStepLR 300ep $+$ calib$/K$ & cosine warm-restart, \code{finetune\_epochs}$/$shot & MultiStepLR (khớp gốc) $+$ calib$/K$ & R56 giống \textperiodcentered{} YOLO khác\\
\hline
Model / data & ResNet/VGG\ldots{} \textperiodcentered{} CIFAR/ImageNet & YOLOv1 \textperiodcentered{} DeepFish & ResNet-56 \textperiodcentered{} CIFAR-10 & Khác (bài toán)\\
\hline
\end{tabularx}}
\normalsize

% ---------------------------------------------------------------------
\section{Nhận xét ngắn (giống/khác gốc \& ảnh hưởng)}

\paragraph{Rebuild ResNet-56 --- faithful nhất.}
Giống gốc gần như hoàn toàn: thuật toán chép trung thành, surgery port, schedule back-loaded per-layer,
hyperparams y hệt (300ep / wd 0.005 / bs 256 / decay 150,225). Khác duy nhất: \textbf{thêm copy BN}
(có cờ \code{-{}-no-bn-copy} để khớp gốc tuyệt đối). $\Rightarrow$ ảnh hưởng pp gốc $\approx 0$; đúng
vai \emph{kiểm chứng thuật toán} (ra $\sim$94.76\% @ 22.4\% là pass).

\paragraph{Rebuild YOLO --- giữ thuật toán gốc, đổi harness $+$ phân bổ (cố ý để fair).}
\begin{itemize}
  \item Giống gốc: thuật toán chọn (chép trung thành), schedule \textbf{back-loaded}, re-rank mỗi shot,
        surgery-first.
  \item Khác gốc (chủ đích): surgery $=$ \code{convert\_to\_lean} (harness mình, có copy BN); phân bổ
        \textbf{uniform} thay per-layer list; finetune cosine.
\end{itemize}
$\Rightarrow$ \textbf{không đụng thuật toán chọn filter}; chỉ đổi phần cắt/chỉnh để đưa vào sân YOLO.

\paragraph{Thuật toán CORING (quan trọng).}
\code{pruning/coring.py} hiện là bản \textbf{chép trung thành loop-based} (KHÔNG vectorized, KHÔNG
chuẩn hoá dấu SVD: \code{get\_u\_svd} trả \code{u[:,:rank]} thẳng như gốc). $\Rightarrow$ kết quả
\textbf{trùng khớp CORING gốc}; đánh đổi là chậm hơn. Cả 2 bản rebuild đều gọi bản này.

% ---------------------------------------------------------------------
\section{Lưu ý cho bảng so chính (YOLO)}
\begin{itemize}
  \item Rebuild YOLO: \code{coring} dùng \textbf{back-loaded}, còn \code{bilevel} dùng \textbf{linear}
        ($r\cdot(it{+}1)/K$, \code{prune.py:187}) $\rightarrow$ \textbf{khác lịch}. Hiện là so
        ``method as-published'', \emph{chưa cô lập importance}. Muốn claim ``thuật toán chọn của tôi
        hơn'' $\rightarrow$ cho cả hai \textbf{cùng lịch} (cùng linear hoặc cùng back-loaded).
  \item Phân bổ \textbf{uniform} (rebuild YOLO) khác \textbf{per-layer list} (gốc/ResNet-56) --- hợp lý
        vì bi-level cũng uniform, nhưng khác cách gốc phân bổ sparsity.
  \item Fair nhất: chạy 2 biến thể CORING --- (a) cùng lịch với bi-level (cô lập importance),
        (b) back-loaded native (đúng paper) --- và báo cáo cả hai.
\end{itemize}

\paragraph{Tóm.} Thuật toán CORING \textbf{giống hệt gốc} ở cả 2 rebuild. ResNet-56 bám gốc mọi mặt
(validate). YOLO giữ thuật toán $+$ schedule back-loaded, chỉ đổi harness / phân bổ / finetune để so
với bi-level --- không làm sai lệch phương pháp gốc.

\end{document}